\section*{The investigation in brief}
Kai Tak's Runway 13 approach joined an instrument-guided descent to a visual right turn over Kowloon. The celebrated $47^\circ$ was a change between published magnetic directions, not the aircraft's bank angle. The engineering problem was to coordinate geometry, speed, descent, visual acquisition and stopping distance within a crowded landscape.

This paper makes that coupling inspectable. A reduced aircraft model constructs an explicit approach, flare, touchdown and full stop. Two retained scenarios have separate computer-assisted continuous-time feasibility certificates. Dated maps, aerial photogrammetry and CAL605 recorder and weather evidence show what would be required to connect those mathematical paths to historical reality.

\begin{table}[H]\centering\small
\caption{Verified model outcomes. Identical illustrative braking: 2 s delay, then $0.15g$. Synthetic entry and corridor; no historical obstacle-clearance claim.}
\begin{tabular}{@{}lrr@{}}\toprule
 & Calm-air witness & Steady-wind witness\\\midrule
Touchdown beyond threshold & 313.49 m & 302.11 m\\
Full-stop station & 2,531.32 m & 2,233.23 m\\
Runway remaining (LDA 2,786 m) & 254.68 m & 552.77 m\\
Certified normalized margin, at least & 0.0397 & 0.15079347\\
Necessary family upper bound & $<0.1559772$ & $\leq0.21674438$\\\bottomrule
\end{tabular}
\end{table}

The wind case imposes a uniform, steady wind from $070^\circ$ true at 21 knots. Its additional 298.09 m of remaining runway is a difference between separately optimized scenarios, not an identified causal wind effect or a demonstrated improvement over a historical procedure. Neither optimum is located by the available bounds.

\subsection*{A route through the paper}
{\small\begin{tabularx}{\textwidth}{@{}rXr@{}}
I & The airport that challenged intuition & p.\,\pageref{part:1}\\
II & Reconstructing a lost airport & p.\,\pageref{part:2}\\
III & The mathematics of the impossible-looking turn & p.\,\pageref{part:3}\\
IV & Can we construct a complete landing? & p.\,\pageref{part:4}\\
V & What happens when wind changes? & p.\,\pageref{part:5}\\
VI & What did the historical evidence reveal? & p.\,\pageref{part:6}\\
VII & What we have proved, and what remains unknown & p.\,\pageref{part:7}\\
VIII & A public research resource for Hong Kong & p.\,\pageref{part:8}\\
 & Mathematical derivations and reproduction & p.\,\pageref{app:dynamics}\\
\end{tabularx}}

\medskip\textbf{The central distinction.} A mathematical certificate verifies a trajectory against declared equations and constraints. An aviation-safety finding requires validated aircraft behaviour, complete relevant obstacles, operational procedures and uncertainty evidence. The current work supplies the former in restricted scenarios, not the latter.
\par\medskip
